Machine-learning weather models are trained one step ahead but used for long autoregressive rollouts, and rollout (multi-step) losses are a common remedy. We isolate this one design choice in a small, fully controlled testbed: a circular 1D CNN with \rhval{main/cnn-k-1/l96/n_params/mean:0} parameters emulates the $36$ slow variables of two-scale Lorenz-96 (fast variables hidden), trained with a $K$-step rollout loss, $K\in\{1,2,4,8\}$, five seeds each, with persistence, climatology and a ridge-fitted linear stencil as reimplemented baselines. All CNNs beat the baselines by a wide margin in ACC lead time (about \rhval{main/cnn-k-1/l96/acc_leadtime/mean:1} time units against \rhval{main/linear-stencil-ridge/l96/acc_leadtime/mean:2} for the linear stencil). Rollout training gives a small, monotone gain at medium range: RMSE at lead 2 time units falls from \rhval{main/cnn-k-1/l96/rmse_t20/mean:3} ($K{=}1$) to \rhval{main/cnn-k-8/l96/rmse_t20/mean:3} ($K{=}8$), i.e.\ $K{=}1$ is \rhval{cmp/main/cnn-k-8/l96/rmse_t20/rel_delta:pct1}\% higher than $K{=}8$ (paired $p=\rhval{cmp/main/cnn-k-8/l96/rmse_t20/paired_p:3}$; $K{=}4$ is only marginal), at the price of a worse one-step RMSE. We find \emph{no} stability or variance benefit to report: every free run of 200 time units survives for every $K$, including one-step training, and the climatological variance ratio stays within about one percent of one for all rollout-length settings. At lead 5 time units the CNN's RMSE exceeds that of climatology although its ACC is still positive. The result is mixed and at small scale: the expected instability of one-step training did not appear in this setting.
